%! TeX program = lualatex
\documentclass[main]{subfiles}

\begin{document}

\title[AutoML: BO]{Important Surrogate Models}
\subtitle{Gaussian Processes and Random Forests}
\author[Bernd Bischl]{Bernd Bischl}
\date{}

\titlebackground*{images/AutoML_HPO_2_2}

%\institute{}
%\date{}

    \maketitle


\begin{frame}{Surrogate Models}
    Desiderata:
    \begin{itemize}
        \item Regression model (there are also classification approaches)
        \item Non-linear local model
        \item Accurate predictions (especially for small sample sizes)
        \item Often: uncertainty estimates
        \item Robust, works often well without human modeler intervention
    \end{itemize}
    Depending on the application:
    \begin{itemize}
        \item Can handle different types of inputs (numerical and categorical)
        \item Can handle dependencies (i.e.\ hierarchical input)
    \end{itemize}
\end{frame}


\begin{frame}{Gaussian Process}
    \small
    Posterior predictive distribution for test point $\conf \in \pcs$ under zero mean:
    $$Y(\conf) \mid \conf, \Dt \sim \normal\left(\surro(\conf), \sh^2(\conf)\right)$$
    with
    $$\surro(\conf) = k(\conf)^\top \bm{K}^{-1} \yv, \quad \sh^2(\conf) = k(\conf, \conf) - k(\conf)^\top \bm{K}^{-1} k(\conf)$$
    \begin{center}
        \includegraphics[width=.4\textwidth]{images/bayesian_loop_sm_normal.png}
    \end{center}
    \footnotesize Note: $\conf$ here denotes the test input. $k(\conf) := (k(\conf, \confI{1}), \ldots, k(\conf, \confI{t}))^\top$, $\yv := (\ysi[1], \ldots, \ysi[t])^\top$. Kernel/Gram matrix $\bm{K} := (k(\confI{i}, \confI{j}))_{i,j}$ for $i, j \in \{1, \ldots, t\}$.
\end{frame}


\begin{frame}{Gaussian Process}
    \small
    Example kernel functions:
    \begin{itemize}
        \item Radial basis function kernel (aka Gauss kernel):
        $$k(\conf, \conf') = \exp\!\left(-\frac{d(\conf, \conf')^2}{2 l^2}\right)$$
        \begin{itemize}
            \item $l$ length scale; $d(\cdot, \cdot)$ Euclidean distance
            \item Infinitely differentiable -- very ``smooth''
        \end{itemize}
        \item Matérn kernels:
        $$k(\conf, \conf') = \frac{1}{\Gamma(\nu)\, 2^{\nu-1}}\left(\frac{\sqrt{2\nu}}{l} d(\conf, \conf')\right)^\nu K_\nu\!\left(\frac{\sqrt{2\nu}}{l} d(\conf, \conf')\right)$$
        \begin{itemize}
            \item $l$ length scale; $K_\nu(\cdot)$ modified Bessel function; $\Gamma(\cdot)$ Gamma function
            \item For $\nu = 3/2$ once differentiable, for $\nu = 5/2$ twice differentiable
            \item Popular choice as a kernel function when using a GP as SM
        \end{itemize}
    \end{itemize}
\end{frame}


\begin{frame}{Gaussian Process}
    \textbf{Pros:}
    \begin{itemize}
        \item Smooth, local, powerful estimator, also for small data
        \item GPs yield well-calibrated uncertainty estimates
        \item The posterior predictive distribution under a GP is normal
    \end{itemize}
    \textbf{Cons:}
    \begin{itemize}
        \item Vanilla GPs scale cubically in the number of data points
        \item Can natively only handle numeric features; mixed inputs / dependencies require special kernels
        \item Not that robust; numerical problems can occur
        \item Can be sensitive to the choice of kernel and hyperparameters
    \end{itemize}
\end{frame}


\begin{frame}{Random Forest}
    \begin{itemize}
        \item Bagging ensemble
        \item Fit $B$ decision trees on bootstrap samples
        \item Feature subsampling
    \end{itemize}
    \vfill
    \begin{center}
        \includegraphics[width=.7\textwidth]{images/random_forests.jpg}
    \end{center}
    ``extratrees'' / random splits:
    \begin{itemize}
        \item Choose split location uniformly at random
        \item Results in a ``smoother'' mean prediction
    \end{itemize}
\end{frame}


\begin{frame}{Random Forest -- Mean and Variance}
    \small
    \begin{itemize}
        \item Let $\surro_b: \pcs \to \R$ be the mean prediction of tree $b$ (mean of all data points in the same node as $\conf$)
        \item Let $\sh_{b}^2: \pcs \to \R$ be the variance prediction (variance of all data points in the same node as $\conf$)
        \item Mean prediction of forest: $\surro: \pcs \to \R, \conf \mapsto \frac{1}{B} \sum_{b=1}^{B} \surro_{b}(\conf)$
        \item Variance prediction of forest:
        $$\conf \mapsto \left(\frac{1}{B} \sum_{b=1}^{B} \sh_{b}^2(\conf) + \surro_{b}(\conf)^{2}\right) - \surro(\conf)^{2}$$
        (law of total variance assuming a mixture of $B$ models)
        \item Alternative variance estimator: (infinitesimal) Jackknife
        \item Variance prediction derived from randomness of individual trees
        \begin{itemize}
            \item Bagging / bootstrap samples
            \item Features sampled at random
            \item (Randomized split locations in the case of ``extratrees'')
        \end{itemize}
    \end{itemize}
\end{frame}


\begin{frame}{Random Forest -- Different Choices}
    \begin{center}
        \includegraphics[width=.95\textwidth]{images/surrogate_1.png}
    \end{center}
\end{frame}


\begin{frame}{Random Forest}
    \textbf{Pros:}
    \begin{itemize}
        \item Cheap(er) to train
        \item Scales well with the number of data points
        \item Scales well with the number of dimensions
        \item Can easily handle hierarchical mixed spaces (via imputation or directly in the tree structure)
        \item Robust
    \end{itemize}
    \textbf{Cons:}
    \begin{itemize}
        \item Suboptimal uncertainty estimates
        \item Not really Bayesian (no real posterior predictive distribution)
        \item Poor extrapolation
    \end{itemize}
\end{frame}


\begin{frame}{Example}
    Minimize the 2D Ackley function using BO\_GP (GP with Matérn 3/2, EI), BO\_RF (standard Random Forest, EI), BO\_RF\_ET (Random Forest with extratrees, EI) or random search:
    \begin{center}
        \includegraphics[width=.6\textwidth]{images/surrogate_2.png}
    \end{center}
    \begin{footnotesize}
        Strong BO\_GP performance. BO\_RF and BO\_RF\_ET not too bad either. BO\_RF\_ET maybe slightly better final performance than BO\_RF.
    \end{footnotesize}
\end{frame}


\begin{frame}{Summary}
\begin{enumerate}
    \item A good BO surrogate provides accurate, uncertainty-aware predictions and handles mixed input types if needed
    \item Gaussian processes are the default choice: well-calibrated uncertainty, but limited scalability and robustness
    \item Random forests scale better and handle mixed/hierarchical spaces, at the cost of weaker uncertainty estimates
\end{enumerate}
\end{frame}

\end{document}
